% =============================================================================
% Block 1 --- Aula 1: measurement (GDP by three approaches, price indices, PPP, beyond GDP).
% Weighted toward Prova 01 error cluster 1 (every approach computed, and they must agree).
\bloco{1}{Measuring aggregates: GDP, price indices, PPP, beyond GDP}{Aula 1 \,\textbullet\, Kurlat chs.~1--2 \,\textbullet\, 10 points}

\questao{(4 points) \textit{Source: Kurlat, P. (2020), \emph{A Course in Modern Macroeconomics}, Exercise 1.3 ``Chained GDP'',
pp.~28--29 --- statement reproduced verbatim, all items (a)--(c).}}

\noindent The country of Fructus produces Apples, Bananas and Cherries. Its production statistics are given below:
\begin{center}\small
\begin{tabular}{@{}l*{6}{r}@{}}
\toprule
 & \multicolumn{2}{c}{\textbf{Apples}} & \multicolumn{2}{c}{\textbf{Bananas}} & \multicolumn{2}{c}{\textbf{Cherries}} \\
\cmidrule(lr){2-3}\cmidrule(lr){4-5}\cmidrule(l){6-7}
 & quantity & price & quantity & price & quantity & price \\
\midrule
2000 & 100 & 40 & 50 & 30 & 200 & 20 \\
2001 & 103 & 39 & 47 & 35 & 206 & 20 \\
2002 & 107 & 39 & 47 & 38 & 213 & 19 \\
2003 & 109 & 39 & 45 & 38 & 215 & 18 \\
2004 & 112 & 40 & 46 & 41 & 224 & 18 \\
2005 & 112 & 39 & 46 & 39 & 227 & 17 \\
2006 & 114 & 40 & 47 & 36 & 237 & 16 \\
2007 & 115 & 40 & 46 & 42 & 249 & 16 \\
2008 & 119 & 40 & 43 & 44 & 255 & 16 \\
2009 & 118 & 42 & 40 & 46 & 262 & 15 \\
2010 & 116 & 43 & 41 & 44 & 268 & 15 \\
2011 & 118 & 42 & 40 & 50 & 280 & 15 \\
2012 & 117 & 42 & 41 & 50 & 289 & 15 \\
2013 & 119 & 43 & 40 & 53 & 290 & 14 \\
2014 & 120 & 44 & 40 & 56 & 297 & 13 \\
2015 & 125 & 45 & 41 & 59 & 308 & 13 \\
2016 & 128 & 45 & 38 & 66 & 320 & 13 \\
2017 & 129 & 46 & 38 & 61 & 330 & 13 \\
2018 & 131 & 47 & 37 & 64 & 344 & 13 \\
2019 & 136 & 47 & 37 & 61 & 353 & 13 \\
2020 & 143 & 47 & 34 & 61 & 361 & 12 \\
\bottomrule
\end{tabular}
\end{center}
\noindent(you can download them as an Excel spreadsheet from the book website)
\begin{enumerate}[label=(\alph*),leftmargin=2.2em,topsep=2pt,itemsep=1pt]
\item Compute a real GDP series at year-2000 prices using base year prices and using the chained method.
\item Plot both series and comment on any differences.
\item What was the average growth rate according to each method?
\end{enumerate}
\noindent\textit{Added to the book's items (not in the book):} use the chained method of Kurlat \S1.2; show the 2000--2001 link in full,
describe the plot of (b) in words, and report both series for 2000, 2005, 2010, 2015 and 2020 and the average growth rates.
\resolespaco
\begin{sol}
\begin{notacao}
\item $i\in\{A,B,C\}$: the three goods of the table (B are bananas, C cherries).
\item $p_{i,t}>0$, $q_{i,t}>0$: price and quantity of good $i$ in year $t\in\{2000,\dots,2020\}$.
\item $Y_t$: real GDP at constant 2000 prices. $Y^{ch}_t$: chained real GDP, expressed in 2000 prices.
\item $g^I_t$, $g^F_t$: growth from $t-1$ to $t$ valued at initial ($t-1$) and at final ($t$) prices. $g_t$: chained growth.
\item $\bar g$: average annual growth rate over 2000--2020.
\end{notacao}
\passo{base-year prices}{Real GDP at constant prices values each year's quantities at the 2000 prices,
$(p_A,p_B,p_C)_{2000}=(40,30,20)$:}
\begin{equation}
Y_t=\textstyle\sum_i p_{i,2000}\,q_{i,t}. \label{b1q1:base}
\end{equation}
Using \eqref{b1q1:base} for 2000 and 2001:
\begin{align*}
Y_{2000}&=40\times100+30\times50+20\times200\\
&=4{,}000+1{,}500+4{,}000=9{,}500,\\
Y_{2001}&=40\times103+30\times47+20\times206\\
&=4{,}120+1{,}410+4{,}120=9{,}650.
\end{align*}
\passo{the 2001 chain link}{Kurlat's chained method measures each year's growth at both years' prices and
takes the geometric average:}
\begin{align}
1+g^I_t&=\frac{\sum_i p_{i,t-1}\,q_{i,t}}{\sum_i p_{i,t-1}\,q_{i,t-1}}, \label{b1q1:gI}\\
1+g^F_t&=\frac{\sum_i p_{i,t}\,q_{i,t}}{\sum_i p_{i,t}\,q_{i,t-1}}, \label{b1q1:gF}\\
1+g_t&=\big[(1+g^I_t)(1+g^F_t)\big]^{1/2}, \label{b1q1:g}\\
Y^{ch}_t&=Y^{ch}_{t-1}\,(1+g_t),\qquad Y^{ch}_{2000}=Y_{2000}. \label{b1q1:chain}
\end{align}
For 2001 the initial prices are the 2000 prices, so \eqref{b1q1:gI} is the ratio of the two sums of the previous step:
\[
1+g^I_{2001}=\frac{9{,}650}{9{,}500}=1.01579 .
\]
At the 2001 prices, $(39,35,20)$:
\begin{align*}
\textstyle\sum_i p_{i,2001}\,q_{i,2001}&=39\times103+35\times47+20\times206\\
&=4{,}017+1{,}645+4{,}120=9{,}782,\\
\textstyle\sum_i p_{i,2001}\,q_{i,2000}&=39\times100+35\times50+20\times200\\
&=3{,}900+1{,}750+4{,}000=9{,}650.
\end{align*}
Substituting these into \eqref{b1q1:gF}:
\[
1+g^F_{2001}=\frac{9{,}782}{9{,}650}=1.01368 .
\]
Using \eqref{b1q1:g}:
\[
1+g_{2001}=(1.01579\times1.01368)^{1/2}=1.01473 ,
\]
a chained growth of $1.473\%$. Using \eqref{b1q1:chain}:
\begin{align*}
Y^{ch}_{2001}&=9{,}500\times(1+g_{2001})\\
&=9{,}500\times(1.014734)\\
&=9{,}639.97 .
\end{align*}
Repeating \eqref{b1q1:gI}--\eqref{b1q1:chain} every year:
\begin{center}\small
\begin{tabular}{@{}lccccc@{}}
\toprule
 & 2000 & 2005 & 2010 & 2015 & 2020 \\
\midrule
At 2000 prices & 9{,}500 & 10{,}400 & 11{,}230 & 12{,}390 & 13{,}960 \\
Chained (2000 prices) & 9{,}500 & 10{,}341.8 & 10{,}950.0 & 11{,}922.8 & 12{,}976.2 \\
\bottomrule
\end{tabular}
\end{center}
\passo{(b) the gap}{The fixed-price series grows faster and the gap widens every year. Cherries expand most ($200\to361$) while
their relative price falls ($20\to12$); bananas shrink ($50\to34$) while theirs rises ($30\to61$). The fixed 2000 prices in
\eqref{b1q1:base} keep weighting cherries at their old, high price, so they overstate growth; \eqref{b1q1:gI}--\eqref{b1q1:gF}
update the weights every year.}
\passo{(c) average growth}{The average annual rate compounds over the 20 years:}
\begin{equation}
1+\bar g=\Big(\frac{Y_{2020}}{Y_{2000}}\Big)^{1/20}. \label{b1q1:avg}
\end{equation}
Applying \eqref{b1q1:avg} to each series of the table:
\begin{align*}
\text{at 2000 prices:}\quad 1+\bar g&=(13{,}960/9{,}500)^{1/20}=1.0194,\\
\text{chained:}\quad 1+\bar g&=(12{,}976.2/9{,}500)^{1/20}=1.0157,
\end{align*}
that is, $1.94\%$ and $1.57\%$ a year.
\intuicao{A base year far in the past weights each good by a relative price that no longer holds; when the goods that grow most are the
ones getting cheaper, that weighting inflates measured real growth.}
\refbib{Kurlat (2020), ch.~1, \S1.2, ``Alternative 3: chained prices'' (pp.~22--24); Exercise~1.5 (p.~29) asks why expanding sectors
tend to get cheaper. The book prints no answer; every figure is recomputed in the check script.}
\rubrica{
\item (a) base-year series 1 pt, chained series with the 2001 link shown 1.5 pts; (b) 0.75 pt; (c) 0.75 pt.
\item Chained growth computed as an arithmetic average of $g^I$ and $g^F$: $-0.25$ pt (the numbers barely change, the method does).
Average growth taken as the arithmetic mean of annual rates: $-0.25$ pt. Series at current prices called ``real'': max.\ 1 pt.
}
\end{sol}

% -----------------------------------------------------------------------------
\questao{(5 points) Discursive. \emph{Three approaches $\times$ price indices $\times$ beyond GDP.} An illustrative Brazilian
economy (BRL billion, 2025) has four producers.
\begin{itemize}[leftmargin=1.4em]
\item Agro firm F grows soybeans with no intermediate inputs. It sells 300 to the food processor P, exports 200 and pays wages of 250.
\item P buys the 300 of soybeans, imports packaging for 100 and buys 40 of online advertising from N. It produces 700 of food:
households buy 650 and 50 stays unsold in P's warehouse. P pays wages of 180.
\item Platform N gives households a free app financed by advertising; its only sale is the 40 of ads to P, and it pays 30 to the app
staff. It also pays 60 in wages to programmers who write software that N will use in its own operations for ten years.
\item The government pays 120 in wages to public school teachers and 80 in Bolsa Fam\'ilia transfers, financed by taxes.
\end{itemize}
Each firm's profit is its revenue plus any output added to its own stocks, minus wages and intermediate inputs.}

\textbf{a)} (2 points) Compute 2025 GDP by the production, income and expenditure approaches, and show that they agree.
\resolespaco
\begin{sol}
\begin{notacao}
\item $j\in\{F,P,N,\text{Gov}\}$: the agro firm, the food processor, the platform and the government (BRL billion, 2025).
\item $VA_j$: value added of producer $j$. $W_j\ge0$: wages it pays. $\Pi_j$: its profit (statement's definition).
\item $C$, $I$, $G$, $X$, $M$: household consumption, investment (fixed capital plus inventories), government purchases
of goods and services, exports and imports.
\end{notacao}
\passo{the four identities used below}{GDP is the value of final goods and services produced; the three approaches must
agree:}
\begin{align}
\text{GDP}&=\textstyle\sum_j VA_j, \label{b1q2:gdpva}\\
VA_j&=\text{gross output}_j-\text{intermediate inputs}_j, \label{b1q2:va}\\
\text{GDP}&=\textstyle\sum_j W_j+\sum_j\Pi_j, \label{b1q2:inc}\\
\text{GDP}&=C+I+G+X-M. \label{b1q2:exp}
\end{align}
The government sells nothing, so its output is valued at cost: its value added is its wage bill.
\passo{production (value added)}{P's gross output is its sales plus the food added to its stocks:
$650+50=700$. Own-account software is valued at cost, 60. Using \eqref{b1q2:va}:}
\begin{align*}
VA_F&=500-0=500,\\
VA_P&=700-300-100-40=260,\\
VA_N&=(40+60)-0=100,\\
VA_{\text{Gov}}&=120.
\end{align*}
Summing with \eqref{b1q2:gdpva}:
\[
\text{GDP}=500+260+100+120=980 .
\]
\passo{income}{Wages, with N paying the app staff and the programmers:}
\[
\textstyle\sum_j W_j=250+180+(30+60)+120=640 .
\]
Profits, by the statement's definition (revenue plus own stocks, minus wages and inputs):
\begin{align*}
\Pi_F&=500-250=250,\\
\Pi_P&=700-300-100-40-180=80,\\
\Pi_N&=100-90=10,
\end{align*}
a total of $340$. Using \eqref{b1q2:inc}:
\[
\text{GDP}=640+340=980 .
\]
\passo{expenditure}{$C=650$ (the app is free: no transaction). $I$ is inventory plus software, which is gross fixed capital
formation: $I=50+60=110$. $G=120$ (the transfer of 80 buys no good or service). $X=200$, $M=100$. Using \eqref{b1q2:exp}:}
\[
\text{GDP}=650+110+120+200-100=980 .
\]
\passo{traps}{Summing gross outputs, $500+700+100+120=1{,}420$, double counts the soybeans and the ads, which P uses up;
\eqref{b1q2:va} nets them out.}
\refbib{Kurlat (2020), ch.~1, \S1.1 (p.~15); Prova~01, Q1(b)--(c) (\texttt{Prova/}\allowbreak\texttt{prova-01-error-analysis.md}).}
\rubrica{
\item 0.7 pt per approach (2 pts in total with agreement). Software placed in $C$, or the transfer in $G$: max.\ 1 pt in the item.
}
\end{sol}

\textbf{b)} (1.5 points) Give the number and the reason in each case. (i) GDP if N books the software as a current expense.
(ii) GDP a student obtains by counting Bolsa Fam\'ilia in $G$. (iii) Gross national income if 30 of F's profit belongs to a foreign
shareholder. (iv) What part of the app's value to households GDP records, and how a welfare measure such as Jones--Klenow's would treat it.
\resolespaco
\begin{sol}
\begin{notacao}
\item Symbols as in (a). GNI: gross national income, the factor income of residents.
\end{notacao}
\passo{(i)}{Expensed, the software is an intermediate input of N. Using \eqref{b1q2:va} and then \eqref{b1q2:exp}, with
$I=50$:}
\begin{align*}
VA_N&=100-60=40,\\
\text{GDP}&=650+50+120+200-100=920 .
\end{align*}
National accounts capitalise own-account software as investment in intellectual property, because it is a durable
productive asset; then GDP $=980$, as in (a).
\passo{(ii)}{$980+80=1{,}060$ is wrong: a transfer moves income from taxpayers to families and pays for no current production,
so it is not in $G$ in \eqref{b1q2:exp}. It enters $C$ only when the families spend it, and that spending is already in the 650.}
\passo{(iii)}{GDP is territorial; GNI counts the income of residents:}
\begin{equation}
\text{GNI}=\text{GDP}-\text{factor income paid abroad}+\text{factor income received from abroad}. \label{b1q2:gni}
\end{equation}
Using \eqref{b1q2:gni}, with nothing received from abroad:
\[
\text{GNI}=980-30+0=950 .
\]
\passo{(iv)}{GDP records only the 40 of ads, as an intermediate input of P that is already inside the price of food. The service
households get at price zero is not a market transaction, so its value is missing from GDP. A consumption-based welfare
measure (Kurlat ch.~2) would count it only if it were valued, e.g.\ by households' willingness to pay or time spent.}
\refbib{Kurlat (2020), ch.~1, \S1.1 and ch.~2, \S2.1--2.2 (pp.~31--33).}
\rubrica{\item 0.375 pt per sub-case, number and reason. Number without the reason: half.}
\end{sol}

\textbf{c)} (1.5 points) In 2026 every quantity repeats, while prices change: food $+10\%$, soybeans $+20\%$ (domestic and
export price), imported packaging $+50\%$, teachers' wages $+5\%$; ads and software unchanged. Compute nominal GDP in 2026, real GDP at
2025 prices and the inflation rate of the GDP deflator. Compare it with a consumer price index (IPCA-type, Laspeyres) built on the
household basket, and explain the gap.
\resolespaco
\begin{sol}
\begin{notacao}
\item Symbols as in (a); subscripts 2025 and 2026 date a value. Quantities $q$ are the same in both years.
\item $P^D_t$: GDP deflator, nominal over real GDP, with $P^D_{2025}=1$. $\pi^D$: its inflation rate.
\item $\text{CPI}_t$: Laspeyres consumer price index on the 2025 household basket $q^H_{2025}$, with $\text{CPI}_{2025}=1$.
\end{notacao}
\passo{nominal GDP}{Each 2025 component of \eqref{b1q2:exp} is multiplied by its price change (software is unchanged):}
\begin{align*}
C&=650\times1.1=715,\\
I&=50\times1.1+60=115,\\
G&=120\times1.05=126,\\
X&=200\times1.2=240,\\
M&=100\times1.5=150.
\end{align*}
Using \eqref{b1q2:exp}:
\begin{equation}
\text{GDP}^{\text{nom}}_{2026}=715+115+126+240-150=1{,}046. \label{b1q2:nom}
\end{equation}
Production check with \eqref{b1q2:va}: F sells $300\times1.2=360$ at home and exports $240$; P's output is $715+55=770$:
\begin{align*}
VA_F&=360+240=600,\\
VA_P&=770-360-150-40=220,\\
VA_N&=100,\qquad VA_{\text{Gov}}=126,
\end{align*}
and \eqref{b1q2:gdpva} gives $600+220+100+126=1{,}046$.
\passo{real GDP and deflator}{At 2025 prices the 2026 quantities give the 2025 value, since nothing moved:}
\begin{equation}
\text{GDP}^{\text{real}}_{2026}=980. \label{b1q2:real}
\end{equation}
The deflator is the ratio of \eqref{b1q2:nom} to \eqref{b1q2:real}:
\[
P^D_{2026}=\frac{1{,}046}{980}=1.0673,\qquad \pi^D=6.73\% .
\]
Real growth is zero, so all of the $6.73\%$ nominal growth is price.
\passo{CPI}{Households buy only food (the app costs zero; the transfer is not a good), so the Laspeyres index is}
\[
\text{CPI}_{2026}=\frac{\sum_i p_{i,2026}\,q^H_{i,2025}}{\sum_i p_{i,2025}\,q^H_{i,2025}}=\frac{650\times1.1}{650}=1.10,
\]
a rise of $10\%$.
\passo{the gap}{The deflator prices domestic production. Splitting the rise from \eqref{b1q2:real} to \eqref{b1q2:nom} by item:
food $700\times0.10=70$, teachers $120\times0.05=6$, exports $200\times0.20=40$, imports $-100\times0.50=-50$
(the soybeans sold to P cancel between F and P):}
\begin{align*}
\text{GDP}^{\text{nom}}_{2026}-\text{GDP}^{\text{real}}_{2026}&=70+6+40-50=66,\\
\pi^D&=66/980=6.73\% .
\end{align*}
Dearer imports \emph{lower} the deflator, because $M$ is subtracted, and exports raise it. The CPI prices what households
consume, imported goods included, with fixed base-year weights.
\refbib{Kurlat (2020), ch.~1, \S1.2 (p.~22).}
\rubrica{
\item Nominal GDP 0.5; deflator 0.5; CPI and the gap 0.5. Imports added rather than subtracted: max.\ 0.5 pt. Arithmetic slip with the right
method: $-0.25$ pt.
}
\end{sol}

% -----------------------------------------------------------------------------
\questao{(1 point) Multiple choice. Brazil's GDP per capita is BRL~60{,}000 (illustrative); the market exchange rate is 5.0 BRL per USD and
the Penn World Table PPP conversion factor is 2.5 BRL per USD; US GDP per capita is USD~85{,}000. Which statement is correct?}
\begin{alternativas}
\item At PPP Brazil's GDP per capita is USD~12{,}000, about 14\% of the US level, since PPP strips out the cheaper Brazilian tradables.
\item At PPP it is USD~24{,}000, about 28\% of the US level, because tradable goods are cheaper in Brazil while the real is undervalued.
\item At PPP it is USD~24{,}000, about 28\% of the US level, because non-tradable services are cheaper in Brazil, so a real buys more there.
\item At PPP it is USD~150{,}000, about 176\% of the US level, because PPP converts reais at the price of a common basket of goods.
\item At market rates it is USD~24{,}000 and at PPP USD~12{,}000, so the market rate overstates Brazilian living standards.
\end{alternativas}
\gabarito{C}
\begin{sol}
\begin{notacao}
\item $y=60{,}000$: GDP per capita in BRL; $e$: a conversion rate in BRL per USD (market $5.0$, PPP $2.5$); $y^{US}=85{,}000$ USD.
\end{notacao}
\passo{numbers}{A value in BRL is converted to USD by dividing by the rate:}
\begin{equation}
y^{USD}=\frac{y}{e}. \label{b1q3:conv}
\end{equation}
Using \eqref{b1q3:conv} at each rate, then dividing by $y^{US}$:
\begin{align*}
\text{market:}\quad y^{USD}&=60{,}000/5.0=12{,}000, & 12{,}000/85{,}000&=14.1\%,\\
\text{PPP:}\quad y^{USD}&=60{,}000/2.5=24{,}000, & 24{,}000/85{,}000&=28.2\%.
\end{align*}
\passo{reason}{The PPP factor is below the market rate because non-tradables (haircuts, rents, services) are cheaper in a
lower-productivity country (Balassa--Samuelson); tradables roughly obey the law of one price. A uses the market number with the wrong
label; B has the right number with the tradables story; D multiplies by the PPP factor ($60{,}000\times2.5$) and keeps a wrong reason
(two mistakes); E swaps the two conversions.}
\refbib{Kurlat (2020), ch.~1, \S1.2 (PPP comparisons, p.~22); Jones (2020), ch.~2.}
\end{sol}
